\documentclass[a4paper,11pt]{article}

\usepackage{psychtools}
\usepackage{hyperref}
\usepackage{booktabs}
\usepackage{microtype}
\usepackage{parskip}
\usepackage{makeidx}

\makeindex

\hypersetup{
  colorlinks=true,
  urlcolor=blue,
  linkcolor=black,
}

\newcommand{\pkg}[1]{\textsf{#1}}
\newcommand{\cmd}[1]{\texttt{\textbackslash #1}}
\newcommand{\opt}[1]{\texttt{#1}}
\newcommand{\meta}[1]{\ensuremath{\langle}\textit{#1}\ensuremath{\rangle}}

% Index helpers
\newcommand{\icmd}[1]{\cmd{#1}\index{#1@\texttt{\textbackslash #1}}}
\newcommand{\iopt}[1]{\opt{#1}\index{#1@\texttt{#1} (option)}}

\title{\pkg{psychtools} --- LaTeX Macros for Psychology Manuscripts\\[0.5ex]
  \large Version 1.0}
\author{Ansgar Endress\\\texttt{ansgar.endress@gmx.net}}
\date{2026/06/26}

\begin{document}
\maketitle
\tableofcontents
\bigskip\hrule\bigskip

%% ---------------------------------------------------------------
\section{Introduction}

\pkg{psychtools} provides a collection of convenience macros for writing
empirical psychology and cognitive science manuscripts in \LaTeX. It addresses
four common needs:

\begin{itemize}
  \item \textbf{Statistical notation} (\S\ref{sec:stats}): compact, italic
        commands for standard test statistics and effect sizes, matching
        APA formatting requirements.
  \item \textbf{Supplementary-materials appendix} (\S\ref{sec:appendix}):
        automatic \texttt{S}-prefixed numbering for pages, sections, figures,
        and tables, with an optional local table of contents.
  \item \textbf{Draft editing aids} (\S\ref{sec:draft}): footnotes and
        section headers that are marked for removal are shown in red during
        writing but can be hidden entirely for the final submission by
        switching a single package option.
  \item \textbf{Citation shortcuts} (\S\ref{sec:cite}): two helpers for
        \pkg{apacite} users.
\end{itemize}

\noindent Load the package in the usual way:
\begin{verbatim}
  \usepackage{psychtools}
  % or with options:
  \usepackage[excludedetails,appendixtoc]{psychtools}
\end{verbatim}

\pkg{psychtools} requires \pkg{tocloft}, \pkg{xcolor}, \pkg{subcaption},
\pkg{xspace}, and \pkg{tikz}, all of which are part of standard \TeX~Live and
Mik\TeX{} distributions.

A self-contained worked example (\texttt{howto.tex} / \texttt{howto.pdf})
is distributed alongside this documentation. It illustrates every command
and option using a simulated serial-position-effect free-recall experiment.

%% ---------------------------------------------------------------
\section{Package Options}
\label{sec:options}

\pkg{psychtools} declares five options, described below. Where a
default is listed the option is active unless you explicitly override
it; options with no default are inactive unless you request them.

\subsection*{\iopt{includedetails} / \iopt{excludedetails}
  \quad (default: \opt{includedetails})}

These two options control the visibility of \emph{draft content} ---
footnotes and section headers that are intended to remind the author
of outstanding work and that should be removed before final
submission.  Draft content is produced by \cmd{footnotewithdetails},
\cmd{draftsection}, \cmd{draftsubsection}, and
\cmd{draftsubsubsection} (see \S\ref{sec:draft}).

With the default \opt{includedetails}, draft content appears in the
output (coloured red unless \opt{nocolorize} is also given).  Switch
to \opt{excludedetails} when preparing the submission copy: all draft
content disappears silently, with no extra blank space or other
artefact.

\subsection*{\iopt{appendixtoc} \quad (default: off)}

When \opt{appendixtoc} is active, \pkg{psychtools} suppresses
supplementary-materials entries from the \emph{main} document table
of contents and instead prints a separate, local TOC immediately
after the supplementary-materials header produced by \cmd{myappendix}
(see \S\ref{sec:appendix}).

Without this option, supplementary sections appear in the main
document TOC normally.

\subsection*{\iopt{nocolorize} \quad (default: off)}

By default, draft content is typeset in red so it is easy to
distinguish from the main text during writing.  Passing
\opt{nocolorize} suppresses the colouring: draft content is shown in
the normal text colour.  This option also affects \cmd{colorize}
(see \S\ref{sec:misc}).

\subsection*{\iopt{nocitecommands} \quad (default: off)}

Disables the \pkg{apacite} citation helpers \cmd{citeeg} and
\cmd{cites} (see \S\ref{sec:cite}).  Use this option if you are not
using \pkg{apacite} and wish to define these command names yourself.

%% ---------------------------------------------------------------
\section{Statistical Notation}
\label{sec:stats}

The commands in Table~\ref{tab:stats} typeset common test statistics and
effect sizes in italic math style, matching APA formatting requirements.
Each command includes a trailing \cmd{xspace} so it can be used freely in
running text without extra braces.

\begin{table}[h]
\centering
\begin{tabular}{lll}
\toprule
Command & Output & Meaning \\
\midrule
\cmd{T}  & \T  & Student's $t$ \\
\cmd{F}  & \F  & $F$ ratio \\
\cmd{Z}  & \Z  & $Z$ score \\
\cmd{p}  & \p  & $p$ value \\
\cmd{M}  & \M  & Mean \\
\cmd{SD} & \SD & Standard deviation \\
\cmd{SE} & \SE & Standard error \\
\cmd{D}  & \D  & Cohen's $d$ \\
\cmd{CI} & \CI & 95\% confidence interval \\
\cmd{et} & \et & Eta-squared \\
\cmd{etp}& \etp& Partial eta-squared \\
\cmd{U}  & \U  & Mann--Whitney $U$ \\
\cmd{W}  & \W  & Wilcoxon $W$ \\
\bottomrule
\end{tabular}
\caption{Statistical notation commands.}
\label{tab:stats}
\end{table}

\paragraph{Example}
\begin{verbatim}
  The effect was significant, \T(38) = 2.45, \p = .019, \D = 0.77.
\end{verbatim}
Produces: The effect was significant, \T(38) = 2.45, \p = .019, \D = 0.77.

\paragraph{Compatibility warning}
\label{sec:compat}
The short command names (\cmd{T}, \cmd{F}, \cmd{Z}, \cmd{p}, \cmd{M},
\cmd{U}, \cmd{W}) may conflict with commands defined by math or font
packages. To minimise the risk, load \pkg{psychtools} \emph{before}
any math-heavy packages. If you encounter a ``command already defined''
error, identify which other package defines the conflicting name and
either reorder the \cmd{usepackage} declarations or stop loading that
package.

%% ---------------------------------------------------------------
\section{Supplementary Materials and Appendix}
\label{sec:appendix}

Many psychology journals require supplementary materials to use a
distinct numbering scheme: pages are labelled \texttt{S1},
\texttt{S2},~\ldots, and figures and tables carry the prefix
\texttt{S} (e.g.\ Figure~S1, Table~S3).  Setting this up manually
requires resetting multiple \LaTeX{} counters and renewing several
formatting commands.  \pkg{psychtools} encapsulates all of this in a
single call.

\subsection{\icmd{myappendix}}

Place \cmd{myappendix} at the point in your document where the
supplementary materials begin. It performs the following actions
automatically:

\begin{itemize}
  \item resets the page counter to 1 and formats it as \texttt{S1},
        \texttt{S2},~\ldots;
  \item resets the section counter to 0 and prefixes section numbers
        with the \emph{short name} derived from \cmd{appendixname}
        (see below);
  \item resets the figure and table counters to 0 and prefixes their
        numbers with \texttt{S};
  \item prints a large bold heading containing \cmd{appendixname};
  \item if the \opt{appendixtoc} option is active, adjusts TOC
        indentation for the longer prefixes and prints a local table
        of contents.
\end{itemize}

\noindent\textbf{Setting the appendix name.}  You should set
\cmd{appendixname} in your document \emph{before} calling
\cmd{myappendix}:

\begin{verbatim}
  \renewcommand{\appendixname}{Supplementary Online Materials}
  \myappendix
\end{verbatim}

If \cmd{appendixname} is not set, the language-dependent default
(e.g.\ ``Appendix'' in English) is used and the prefix will be
\texttt{S}. The following values produce a two-letter prefix via
\cmd{appendixshortname}:

\begin{center}
\begin{tabular}{ll}
\toprule
\cmd{appendixname} & Prefix \\
\midrule
\texttt{Supplementary Information}      & \texttt{SI} \\
\texttt{Supplementary Materials}        & \texttt{SM} \\
\texttt{Supplementary Online Materials} & \texttt{SM} \\
anything else                           & \texttt{S}  \\
\bottomrule
\end{tabular}
\end{center}

\subsection{\icmd{appsection}\{\meta{title}\}}

Use \cmd{appsection} for all sections inside the supplementary
materials.  It behaves like a regular \cmd{section} with one
important difference: from the \emph{second} section onwards it
automatically inserts a \cmd{clearpage} before the new section
heading, so that each supplementary section starts on a fresh page.
The first section is not preceded by a page break (since
\cmd{myappendix} already ends with \cmd{clearpage}).

This differs from a plain \cmd{section}, which never inserts an
automatic page break.  Using \cmd{section} inside the supplementary
materials works, but you would need to insert \cmd{clearpage}
manually between sections.

\begin{verbatim}
  \renewcommand{\appendixname}{Supplementary Online Materials}
  \myappendix
  \appsection{Stimuli}          % first section — no preceding page break
  ...
  \appsection{Additional Analyses}  % second section — preceded by \clearpage
  ...
\end{verbatim}

\subsection{\icmd{appsectionFirst}\{\meta{title}\} (deprecated)}

An earlier version of \pkg{psychtools} required the \emph{first}
appendix section to be introduced with \cmd{appsectionFirst} rather
than \cmd{appsection}, because the original \cmd{appsection}
unconditionally inserted a \cmd{clearpage}.  The current
implementation of \cmd{appsection} checks the section counter and
skips the page break when it is zero, making \cmd{appsectionFirst}
a no-op alias.  It is retained for backward compatibility; new
documents should use \cmd{appsection} throughout.

\subsection{\icmd{appendixshortname}}

\cmd{appendixshortname} expands to the short prefix (\texttt{SI},
\texttt{SM}, or \texttt{S}) used for section numbers.  It is called
internally by \cmd{myappendix} but is available if you need the
prefix elsewhere in your document.

%% ---------------------------------------------------------------
\section{Draft Editing Aids}
\label{sec:draft}

During manuscript preparation it is common to keep reminders about
content that still needs to be added or revised --- for example, a
section that is present only as a placeholder, or a footnote
explaining a methodological choice that will be cut before
submission.  The commands in this section allow you to include such
material in the source file while making it visually prominent
(coloured red) and easy to remove: switching the \opt{excludedetails}
package option hides all draft content without leaving blank space or
other artefacts.

\subsection{Draft section headings}

\begin{center}
\begin{tabular}{lp{9cm}}
\toprule
Command & Effect \\
\midrule
\icmd{draftsection}\{\meta{t}\} &
  Typesets a \cmd{section} heading with the prefix
  \colorize{(Header will be removed)} in red.
  The heading is invisible with \opt{excludedetails}. \\[4pt]
\icmd{draftsubsection}\{\meta{t}\} &
  Same for a \cmd{subsection} heading. \\[4pt]
\icmd{draftsubsubsection}\{\meta{t}\} &
  Same for a \cmd{subsubsection} heading. \\
\bottomrule
\end{tabular}
\end{center}

When the \opt{appendixtoc} option is active, draft headings are
automatically excluded from the supplementary-materials TOC.

\subsection{\icmd{footnotewithdetails}\{\meta{text}\}}

Inserts a footnote with the prefix \colorize{(Footnote will be
removed)} in red, followed by \meta{text}.  The footnote is
invisible with \opt{excludedetails}.  Use this for methodological
notes or explanations that are helpful during review but should be
cut from the final paper.

\subsection{\icmd{colorize}\{\meta{text}\}}

The primary use of \cmd{colorize} is to mark passages that have been
revised in response to reviewer comments, so that the editor and
reviewers can locate the changes easily when reading the revised
manuscript.  \meta{text} is typeset in red; with the \opt{nocolorize}
package option it is typeset in the normal text colour instead (useful
when submitting a clean version without visible markup).

\cmd{colorize} is also used internally by the draft-heading and
\cmd{footnotewithdetails} commands.

%% ---------------------------------------------------------------
\section{Citation Helpers (apacite)}
\label{sec:cite}

These commands are active when \pkg{apacite} is loaded and the
\opt{nocitecommands} option is not given.  If \pkg{apacite} is not
loaded, basic fallbacks using \cmd{cite} and
\cmd{citeauthor}/\cmd{citeyear} are provided, and a package warning
is issued.

\begin{center}
\begin{tabular}{lll}
\toprule
Command & Example output & Use case \\
\midrule
\icmd{citeeg}\{\meta{key}\} &
  (e.g., Smith 2020) &
  Cite an example reference \\
\icmd{cites}\{\meta{key}\} &
  Smith's (2020) &
  Possessive citation in running text \\
\bottomrule
\end{tabular}
\end{center}

%% ---------------------------------------------------------------
\section{Miscellaneous}
\label{sec:misc}

\subsection{\icmd{includesubgraphics}\{\meta{label}\}\{\meta{filename}\}}

Inserts a subfigure (via the \pkg{subcaption} package) at 90\% of
the text width, with a bold large label (e.g.\ \textbf{A},
\textbf{B}) but \emph{without} a sub-caption.  Use this when you
want panel labels in a multi-panel figure but prefer a single shared
caption below the whole figure rather than individual captions per
panel.

\begin{verbatim}
  \begin{figure}
    \includesubgraphics{A}{panel_a.pdf}
    \includesubgraphics{B}{panel_b.pdf}
    \caption{Results. Panel A shows ...; Panel B shows ...}
  \end{figure}
\end{verbatim}

\subsection{\icmd{protectedemail}\{\meta{address}\}}

Renders an email address in \texttt{typewriter} font inside a
\pkg{tikz} node.  Wrapping the address in a TikZ node rather than
plain text prevents some simple web crawlers from harvesting it as a
mailto target.

\begin{verbatim}
  \protectedemail{ansgar.endress@gmx.net}
\end{verbatim}

%% ---------------------------------------------------------------
\section{Backward Compatibility}

An \texttt{ansgar.sty} shim is distributed alongside
\pkg{psychtools}.  It contains a single line:
\begin{verbatim}
  \RequirePackage{psychtools}
\end{verbatim}
Documents that used \cmd{usepackage}\{ansgar\} continue to work
without modification.

%% ---------------------------------------------------------------
\section{License}

\pkg{psychtools} is distributed under the
\href{https://www.latex-project.org/lppl.txt}{LaTeX Project Public
License}, version~1.3c or later.

%% ---------------------------------------------------------------
\appendix
\section{Command and Option Reference}
\label{sec:ref}

\subsection*{Package options}
\begin{center}
\begin{tabular}{lll}
\toprule
Option & Default & See \\
\midrule
\opt{includedetails}  & on  & \S\ref{sec:options} \\
\opt{excludedetails}  & off & \S\ref{sec:options} \\
\opt{appendixtoc}     & off & \S\ref{sec:options} \\
\opt{nocolorize}      & off & \S\ref{sec:options} \\
\opt{nocitecommands}  & off & \S\ref{sec:options} \\
\bottomrule
\end{tabular}
\end{center}

\subsection*{Commands}
\begin{center}
\begin{tabular}{ll}
\toprule
Command & See \\
\midrule
\cmd{T}, \cmd{F}, \cmd{Z}, \cmd{p}, \cmd{M} & \S\ref{sec:stats} \\
\cmd{SD}, \cmd{SE}, \cmd{D}, \cmd{CI}        & \S\ref{sec:stats} \\
\cmd{et}, \cmd{etp}, \cmd{U}, \cmd{W}        & \S\ref{sec:stats} \\
\cmd{myappendix}           & \S\ref{sec:appendix} \\
\cmd{appsection}           & \S\ref{sec:appendix} \\
\cmd{appsectionFirst}      & \S\ref{sec:appendix} \\
\cmd{appendixshortname}    & \S\ref{sec:appendix} \\
\cmd{draftsection}         & \S\ref{sec:draft} \\
\cmd{draftsubsection}      & \S\ref{sec:draft} \\
\cmd{draftsubsubsection}   & \S\ref{sec:draft} \\
\cmd{footnotewithdetails}  & \S\ref{sec:draft} \\
\cmd{colorize}             & \S\ref{sec:draft} \\
\cmd{citeeg}               & \S\ref{sec:cite} \\
\cmd{cites}                & \S\ref{sec:cite} \\
\cmd{includesubgraphics}   & \S\ref{sec:misc} \\
\cmd{protectedemail}       & \S\ref{sec:misc} \\
\bottomrule
\end{tabular}
\end{center}

\printindex

\end{document}
